\scanpage{194}

le fait jusqu'à $u_{\ell-1}$ \deleted{\unsure{$(\ell\leq k)$}}\inserted{$(0\leq\ell\leq k-1)$} prouvons le pour $u_\ell$.

\deleted{Soit} Soit $x$ un élément de $\operatorname{Ker}u_\ell$, \deleted{\unsure{on va montrer}, \unsure{prouvons} qu'il appartient à tout $E_{\ell'}$ $(\ell'\leq\ell)$ (identifié à un sous-espace de $E_\ell$), \unsure{par|$+0$} \unsure{récurrence} et on aura fini. \unsure{On prend}} \inserted{s'il n'était} pas nul il existerait un \deleted{premier} \deleted{\unsure{dernier}} \inserted{\unsure{plus haut|plus grand}} $\ell'$ \deleted{\unsure{$(0\leq\ell'<\ell)$|$(0\leq\ell'\leq\ell)$}}\inserted{$(-1\leq\ell'<\ell)$}

tel que $x\notin E_{\ell'}$ -- $E_{\ell'}$ identifié à un sous-espace de $E_\ell$. Soit $D$ la droite engendrée par $x$, \deleted{soit} $(D\subset E_{\ell'+1})$ et soit $\deleted{\unsure{E|E_j}}F$ un hyperplan de $E_\ell$ \deleted{\unsure{de}} contenant $E_{\ell'}$ et suppl. de $D$ ; soit pour $0\leq j\leq k$,

\margin{$E_{\ell'}\quad E_{\ell'+1}\ldots E_\ell$}

\[\deleted{E'_j\unsure{=}\left\{\begin{array}{ll}E_j&\text{si }j\leq\ell'\\D&\text{si}\end{array}\right.}\]

\[E'_j=E_j\deleted{\unsure{\cap F}}\quad\text{si}\quad j\leq\ell'\quad\text{ou}\quad\deleted{\unsure{E_j=E_j|E'_j=E_j}}\inserted{j>\ell},\quad\deleted{\unsure{j>\ell}}\inserted{E'_j=E_j\cap F\quad\text{si }\ell'<j\leq\ell}\]

\[ E''_j=0\quad\text{si}\quad j\leq\ell'\ \text{ou}\ j>\ell,\qquad E''_j=D\quad\text{si }\ell'<j\leq\ell \]

Alors $E_j=E'_j\oplus E''_j$ pour tout $j$, \deleted{\unsure{les $E'_j$ et}} \inserted{et $\sum E'_j$ et} $\sum E''_j$ sont stables sous $\deleted{\unsure{D}}\unsure{\mathfrak g|\mathscr G}$ et $\ne0$, donc $\sum E_i$ ne serait pas indécomposable.

\uline{Lemme 2} Les $u_i$ $(0\leq i\leq k-1)$ sont \uline{surjectifs}. Se voit soit directement de façon analogue, soit en \deleted{\unsure{tenant}} appliquant le lemme 1 à la repr. contragrédiente.

% archive sheet number: 193
